\chapter{Overview}\label{ch:introduction}

\section{Introduction}

Small Internet-of-Things systems increasingly use 16- and 32-bit microcontrollers, while the requirements on development time and reliability continue to grow.
A lightweight system is therefore needed to build simple and efficient IoT applications.
At the same time, lightweight virtualisation is becoming more important, so a small operating system is also useful as a guest under another operating system.
Because such systems usually have no memory protection unit, formal verification is needed to establish their reliability.

RMP is a simple, efficient, and easy-to-use fully preemptive real-time operating system (RTOS).
It provides the features expected of a typical RTOS kernel: a fully preemptive static-priority round-robin scheduler, simple memory management, concise communication mechanisms, and hardware-specific optimisations.
RMP is designed to run efficiently on a microcontroller with only 2\,KiB of ROM and 1\,KiB of RAM.

This manual describes the RMP kernel API from the user's point of view.
It first reviews a few concepts of small real-time operating systems, then describes the features and APIs of RMP chapter by chapter.

\subsection{Design goals and specifications}

The goal of RMP is to create a simple, easy-to-use, open-source RTOS kernel with the best real-time behaviour and reliability in its class, and with average execution efficiency comparable to mainstream kernels.

\subsection{Copyright and license}

RMP is released under \emph{The Unlicense}.
Special cases are handled by special provisions, and the exact terms differ for each possible case.\footnote{For example, security equipment and medical equipment.}

\subsection{Terms and definitions}

The basic terms used in this manual are defined as follows.

\subsubsection{Operating system}

The basic software that runs on a device and manages the processor, memory, and other hardware resources.

\subsubsection{Thread}

A schedulable entity that has its own execution stack and instruction stream.
A process may contain several threads that share the process address space.

\subsubsection{Static allocation}

A resource-allocation scheme in which the allocation is decided at compile time.

\subsubsection{Dynamic allocation}

A resource-allocation scheme in which the allocation can change while the system is running.

\subsubsection{Soft real-time}

A real-time guarantee under which most operations complete within their deadlines, but a small fraction may occasionally finish late.

\subsubsection{Hard real-time}

A real-time guarantee under which every operation must complete within its deadline.

\subsubsection{Constant real-time}

A real-time guarantee under which every operation is O(1) in the user input and the system configuration, and every execution completes within a practically meaningful constant deadline.
This is the strongest form of real-time guarantee.

\subsubsection{Value-constant real-time}

A real-time guarantee under which every operation and response is O(1) while a particular value is unchanged, and every execution completes within a practically meaningful constant deadline.

\subsection{Major reference systems}

The scheduler and API are partly modelled on FreeRTOS (RT Engineering Ltd.
/Amazon).
The formal-proof approach is informed by seL4 (2016 Data61/CSIRO).
The references and source material for the other chapters are listed at the end of each chapter.

\section{Foreword}

An operating system is the basic software that runs on a device and manages the processor, memory, and other hardware resources.
For a real-time operating system, every operation must be correct and timely, and its execution time must be predictable.
There are two broad classes of real-time systems: soft real-time systems, where most responses meet their deadlines, and hard real-time systems, where every response must meet its deadline.
In practice, few real-time applications or operating systems are purely soft or purely hard; most combine both.
A common example is an LCD human-machine interface, which is soft real-time, combined with a motor-control loop, which is hard real-time.

RMP is a basic real-time system: the smallest system that exhibits the essential properties of a real-time system.
Such systems generally run on high-end 8/16-bit machines and low-end 32-bit machines and need a system timer.
They do not distinguish user mode from kernel mode, but an MMU or MPU can be configured to protect a region of memory.
They need a small amount of architecture-specific assembly for context-stack switching.
Porting to a new architecture always involves this assembly, and often also the system timer, stack switching, interrupt management, and coprocessor management.
Such systems may use a custom linker script or may not; a custom script is usually required when memory protection is used.

In such a system a task is represented by a thread.
A task function may be reentrant.
Each task has its own execution stack.
Task code and kernel code may be compiled together or separately.
A task calls a system function with an ordinary function call; there is no system-call instruction or software interrupt gate.

Such systems have priorities, and generally implement preemption between different priorities and round-robin scheduling at the same priority.
They have a basic memory-management scheme, usually based on a slab or buddy allocator.

Interrupts may be completely transparent to the operating system, in which case the operating system need not know that an interrupt occurred.
If a context switch is needed in an interrupt, the stack-switch assembly must be inserted into that interrupt, and the interrupt entry and exit must be written in assembly.

Typical systems in this class include RMProkaron, RT-Thread, FreeRTOS, uC/OS, Salvo, and ChibiOS.
Complete formal verification of such a system is usually feasible because the kernel is very small.
Traditional software engineering can systematically test functionality and show that defects are unlikely; formal verification can prove, from a specification or a set of logical rules, that the system has certain properties, which amounts to proving that under the stated conditions the system cannot have those defects.

\section{Performance and specifications of RTOSes}

Hundreds of RTOSes exist, and hobbyist kernels are countless.
Their performance varies widely.
The metrics below are directly comparable only when the processor architecture, the compiler, and the compiler options are the same.
With a different architecture, compiler, or option set, the numbers are indicative but not comparable.
A recommended method is to use an industry-standard ARM or MIPS processor with GCC \texttt{-O2}, or the Chronos simulator.
The system load can also affect the values, so they are comparable only at the same load level.

\subsection{Kernel size}

Kernel size is an important metric for an RTOS, because RTOSes are deployed on devices with very limited memory.
It is measured in two parts: the read-only section, which contains code and read-only data, and the data section, which contains read-write data.
On a flash-based microcontroller, the read-only section consumes flash and the data section consumes SRAM [1].

An RTOS is highly configurable, so its size is not fixed; it depends on the selected configuration.
The metric should therefore be reported for the smallest, a typical, and the largest configuration.

Kernel size is easy to obtain: compile the kernel, then use a binary viewer to inspect the size of each section.\footnote{For example, \texttt{objdump}.}

\subsection{Execution time, worst-case execution time, and jitter}

Execution time is the time taken by an RTOS system call.
Worst-case execution time (WCET) is the maximum execution time under the least favourable conditions.
An RTOS usually reaches its WCET when it executes the longest system call and suffers many cache misses along the way.
System calls normally disable interrupts, so the WCET is often the longest interval with interrupts disabled; its effect on real-time behaviour is therefore very large.

There are two kinds of WCET: the WCET of a kernel system call and the WCET of inter-thread synchronisation.
To measure the first, record a timestamp \(T_s\) before the call and \(T_e\) after it.
Call the timer twice more to obtain \(T_{ts}\) and \(T_{te}\), so the timer overhead is \(T_{te}-T_{ts}\).
The execution time is \(T_e-T_s-(T_{te}-T_{ts})\).
The WCET is the largest such value over all system-call tests.

To measure the second, record \(T_s\) at the sender and \(T_e\) at the receiver.
The timer overhead is measured as above, and the execution time is again \(T_e-T_s-(T_{te}-T_{ts})\).
The WCET is the largest value over all communication tests.

Jitter is also important.
Repeated measurements of the execution or communication time produce a distribution.
Its mean is the average execution time, and its standard deviation\footnote{Sometimes the range is used instead.} is the execution-time jitter.

For an RTOS, execution time, WCET, and jitter are all better when smaller.
Execution time can be divided into the following categories.

\subsubsection{Thread switch time}

Thread switch time is the time to switch from one thread to another.
It can be measured with the \(T_e-T_s\) method or, in a two-thread loop, as half the difference between two successive \(T_s\) values.
There are two cases: a switch between threads of the same priority, and a low-priority thread waking a high-priority thread.

In the first case, assume the two threads have the same priority and that thread A has just been switched in from thread B.

\begin{table}[htbp]
    \centering
    \caption{Thread switch, first case.}
    \small
    \begin{tabular}{@{}p{0.44\textwidth}p{0.44\textwidth}@{}}
        \toprule
        Thread A & Thread B\\
        \midrule
        forever \{ & forever \{ \\ \quad record \(T_s\); & \quad record \(T_e\); \\ \quad switch to thread B; & \quad switch to thread A; \\ \} & \}\\
        \bottomrule
    \end{tabular}
\end{table}

In the second case, thread B has the higher priority and thread A has just been switched in from B.

\begin{table}[htbp]
    \centering
    \caption{Thread switch, second case.}
    \small
    \begin{tabular}{@{}p{0.44\textwidth}p{0.44\textwidth}@{}}
        \toprule
        Thread A & Thread B\\
        \midrule
        forever \{ & forever \{ \\ \quad record \(T_s\); & \quad record \(T_e\); \\ \quad wake B; & \quad sleep; \\ \} & \}\\
        \bottomrule
    \end{tabular}
\end{table}

\subsubsection{Asynchronous communication time}

Asynchronous communication time is the time to send and receive an asynchronous signal between threads.
Assume thread B is already blocked at the receiver, thread A sends, and B has higher priority than A.

\begin{table}[htbp]
    \centering
    \caption{Asynchronous communication between threads.}
    \small
    \begin{tabular}{@{}p{0.44\textwidth}p{0.44\textwidth}@{}}
        \toprule
        Thread A & Thread B\\
        \midrule
        forever \{ & forever \{ \\ \quad record \(T_s\); & \quad record \(T_e\); \\ \quad send to B; & \quad receive from own mailbox; \\ \} & \}\\
        \bottomrule
    \end{tabular}
\end{table}

\subsection{Interrupt response time, worst-case response time, and jitter}

Interrupt response time is the time from an interrupt request to the RTOS invoking the handler thread.
Worst-case interrupt response time is its maximum under the least favourable conditions, typically with many cache and TLB misses during interrupt handling.
This is the most important RTOS metric; one could say that every RTOS design decision revolves around it.
It is the most direct measure of the RTOS's response to an external stimulus.

To measure the worst case, record \(T_s\) in the first assembly instruction of the interrupt vector\footnote{It cannot wait until a C function is entered, because register and stack maintenance are also part of the interrupt response time.} and \(T_e\) at the first line of the handler thread.
The timer overhead is measured as above, and the response time is \(T_e-T_s-(T_{te}-T_{ts})\).
The worst case is the largest value over all interrupt tests.
Jitter is the standard deviation\footnote{Sometimes the range is used instead.} of the distribution.

Smaller response time, worst-case response time, and jitter are better.
The measurement is usually arranged as follows.

\begin{table}[htbp]
    \centering
    \caption{Interrupt response time.}
    \small
    \begin{tabular}{@{}p{0.44\textwidth}p{0.44\textwidth}@{}}
        \toprule
        Kernel & Thread A\\
        \midrule
        hardware interrupt vector \{ & forever \{ \\ \quad record \(T_s\); & \quad record \(T_e\); \\ \quad signal endpoint P; & \quad receive from endpoint P; \\ \} & \}\\
        \bottomrule
    \end{tabular}
\end{table}

\subsection{Realistic interrupt response time, worst case, and jitter}

Realistic interrupt response time is the time from an external signal input to the corresponding I/O output.
Its worst case includes the factors that affect the ordinary worst case plus the intrinsic latency of the CPU and I/O hardware.

Measuring it requires external hardware.
For example, connect one FPGA pin to a CPU input pin and another to a CPU output pin.
The FPGA emits a signal and starts its high-resolution timer; when it receives the output signal, it stops the timer.
The elapsed time is the realistic response time.
The worst case is the largest value over all tests.

For a hardware/software system, smaller realistic response time, worst case, and jitter are better.
The realistic worst case is approximately the worst-case execution time plus the worst-case interrupt response time plus the intrinsic CPU/IO delay.
If an I/O input arrives just as a system call begins, the hardware vector cannot run until the call finishes; the output is produced only after that.
The measurement is arranged as follows.

\begin{table}[htbp]
    \centering
    \caption{Realistic interrupt response time.}
    \small
    \begin{tabular}{@{}p{0.44\textwidth}p{0.44\textwidth}@{}}
        \toprule
        FPGA or oscilloscope & System under test\\
        \midrule
        forever \{ & forever \{ \\ \quad emit signal and start timer; & \quad receive signal from I/O; \\ \quad receive signal; & \quad minimal internal processing; \\ \quad stop timer; & \quad output signal to I/O; \\ \} & \}\\
        \bottomrule
    \end{tabular}
\end{table}

\section{RMP system call format}

A system call is the only way to use a service provided by the system.
In a library-structured operating system such as RMP, a system call is an ordinary function call.
Table~\ref{tab:syscalls} lists all RMP system-call interfaces.

\begin{table}[htbp]
    \centering
    \caption{RMP system calls.}
    \label{tab:syscalls}
    \small
    \begin{tabularx}{\textwidth}{@{}>{\ttfamily}p{4.2cm}rX@{}}
        \toprule
        \tablehead{System call} & \tablehead{No.} & \tablehead{Meaning}\\
        \midrule
        RMP\_Thd\_Yield & 0 & Yield the current thread's processor.\\
        RMP\_Thd\_Crt & 1 & Create a thread.\\
        RMP\_Thd\_Del & 2 & Delete a thread.\\
        RMP\_Thd\_Set & 3 & Set a thread's priority and time slice.\\
        RMP\_Thd\_Suspend & 4 & Suspend a thread.\\
        RMP\_Thd\_Resume & 5 & Resume a suspended thread.\\
        RMP\_Thd\_Delay & 6 & Delay for a period.\\
        RMP\_Thd\_Cancel & 7 & Cancel a thread's delay.\\
        RMP\_Thd\_Snd & 8 & Send a value to a thread's mailbox.\\
        RMP\_Thd\_Snd\_ISR & 9 & Send to a mailbox from an interrupt.\\
        RMP\_Thd\_Rcv & 10 & Receive from the calling thread's mailbox.\\
        RMP\_Sem\_Crt & 11 & Create a counting semaphore.\\
        RMP\_Sem\_Del & 12 & Delete a counting semaphore.\\
        RMP\_Sem\_Pend & 13 & Try to acquire a semaphore.\\
        RMP\_Sem\_Abort & 14 & Abort a thread's semaphore wait.\\
        RMP\_Sem\_Post & 15 & Post to a semaphore.\\
        RMP\_Sem\_Post\_ISR & 16 & Post to a semaphore from an interrupt.\\
        \bottomrule
    \end{tabularx}
\end{table}

\section{Chapter bibliography}

\nocite{anh2009,kar1990,boger2013}
\printbibliography[heading=none]
